\documentclass[a4paper,12pt,twocolumn]{article}

\usepackage[landscape,top=1cm,bottom=1.4cm,left=1cm,right=1cm,columnsep=1cm]{geometry}
\usepackage[fleqn]{amsmath}
\usepackage{amssymb}
\usepackage{fontspec}
\usepackage{enumitem}
\usepackage{xeCJK}

\setlength{\mathindent}{1.5em}

% Set fonts
\setmainfont{Times New Roman}
\setsansfont{Arial}
\setmonofont{JetBrains Mono}

% Chinese font support
\setCJKmainfont{Iansui}[
  BoldFont=Iansui,
  AutoFakeBold=3.17
]

\pagestyle{plain}
\setlength{\columnseprule}{0.2pt}

\newcommand{\examdate}{2026/03/11}
\newlength{\problemnumberwidth}
\settowidth{\problemnumberwidth}{\textbf{7.1.56.}\ }
\newcommand{\problem}[2]{%
  \par\noindent
  \hangindent=\problemnumberwidth
  \hangafter=1
  \makebox[\problemnumberwidth][l]{\textbf{#1.}}#2\par
  \vspace{1.4em}%
}

\begin{document}

\twocolumn[
{\centering
\Large\bfseries Calculus Problems\par
\vspace{0.5em}
\noindent\hfill\normalsize\normalfont \examdate\par
\vspace{0.8em}
}
]

\problem{49}{Determine whether the series is convergent or divergent. If it is convergent, find its sum:
\[
\sum_{n=1}^{\infty}\left(\frac{1}{e^n}+\frac{1}{n(n+1)}\right)
\]}

\problem{50}{Determine whether the series is convergent or divergent. If it is convergent, find its sum:
\[
\sum_{n=1}^{\infty}\frac{e^n}{n^2}
\]}

\problem{77}{Find the value of \(c\) if
\[
\sum_{n=2}^{\infty}(1+c)^{-n}=2
\]}

\problem{8}{Use the Integral Test to determine whether the series is convergent or divergent:
\[
\sum_{n=1}^{\infty} n^2 e^{-n^3}
\]}

\problem{10}{Use the Integral Test to determine whether the series is convergent or divergent:
\[
\sum_{n=1}^{\infty}\frac{\tan^{-1} n}{1+n^2}
\]}

\problem{32}{Find the values of \(p\) for which the series is convergent:
\[
\sum_{n=3}^{\infty} \frac{1}{n \ln n \,[\ln(\ln n)]^p}
\]}

\problem{34}{Find the values of \(p\) for which the series is convergent:
\[
\sum_{n=1}^{\infty} \frac{\ln n}{n^p}
\]}

\problem{46}{For what values of \(p\) does the series
\[
\sum_{n=2}^{\infty} \frac{1}{n^p \ln n}
\]}

\problem{41}{Approximate the sum of the series correct to four decimal places:
\[
\sum_{n=1}^{\infty} \frac{(-1)^n}{(2n)!}
\]}

\problem{43}{Approximate the sum of the series correct to four decimal places:
\[
\sum_{n=1}^{\infty} (-1)^n n e^{-2n}
\]}

\problem{40}{The Bessel function of order 1 is defined by
\[
J_1(x)=\sum_{n=0}^{\infty}\frac{(-1)^n x^{2n+1}}{n!(n+1)!2^{2n+1}}
\]
\begin{enumerate}[label=(\alph*), leftmargin=2em]
\item Find the domain of \(J_1\).
\item Show that \(J_1\) satisfies the differential equation
\[
x^2 J_1''(x)+x J_1'(x)+(x^2-1)J_1(x)=0
\]
\item Show that \(J_0'(x)=-J_1(x)\).
\end{enumerate}}

\problem{49}{Use the power series for \(\tan^{-1}x\) to prove the following expression for \(\pi\) as the sum of an infinite series:
\[
\pi = 2\sqrt{3}\sum_{n=0}^{\infty}\frac{(-1)^n}{(2n+1)3^n}
\]}

\problem{38}{The period of a pendulum with length \(L\) that makes a maximum angle \(\theta_0\) with the vertical is
\[
T = 4\sqrt{\frac{L}{g}} \int_0^{\pi/2} \frac{dx}{\sqrt{1-k^2\sin^2 x}}
\]
where \(k=\sin\left(\frac{1}{2}\theta_0\right)\) and \(g\) is the acceleration due to gravity.
(In Exercise 7.7.42 we approximated this integral using Simpson's Rule.)
\begin{enumerate}[label=(\alph*), leftmargin=2em]
\item Expand the integrand as a binomial series and use the result of Exercise 7.1.56 to show that
\[
T = 2\pi\sqrt{\frac{L}{g}}
\left[
1+\frac{1^2}{2^2}k^2+\frac{1^2 3^2}{2^2 4^2}k^4+\frac{1^2 3^2 5^2}{2^2 4^2 6^2}k^6+\cdots
\right]
\]
If \(\theta_0\) is not too large, the approximation
\[
T \approx 2\pi\sqrt{\frac{L}{g}}
\]
obtained by using only the first term in the series, is often used. A better approximation is obtained by using two terms:
\[
T \approx 2\pi\sqrt{\frac{L}{g}}\left(1+\frac{1}{4}k^2\right)
\]
\item Notice that all the terms in the series after the first one have coefficients that are at most \(\frac{1}{4}\). Use this fact to compare this series with a geometric series and show that
\[
2\pi\sqrt{\frac{L}{g}}\left(1+\frac{1}{4}k^2\right)
\leq
T
\leq
2\pi\sqrt{\frac{L}{g}}\frac{4-3k^2}{4-4k^2}
\]
\item Use the inequalities in part (b) to estimate the period of a pendulum with \(L=1\) meter and \(\theta_0=10^\circ\). How does it compare with the estimate
\[
T \approx 2\pi\sqrt{\frac{L}{g}} \, ?
\]
What if \(\theta_0=42^\circ\)?
\end{enumerate}}

\problem{7.1.56}{Prove that, for even powers of sine,
\[
\int_0^{\pi/2}\sin^{2n}x\,dx
=
\frac{1\cdot 3\cdot 5\cdots(2n-1)}{2\cdot 4\cdot 6\cdots 2n}\cdot\frac{\pi}{2}
\]}

\problem{3}{Find an equation of sphere with center \((-3,2,5)\) and radius \(4\). What is the intersection of this sphere with the \(yz\)-plane?
\((6\text{ point})\)}

\end{document}
